% STATUS: TODO
% PURPOSE: First Lyapunov coefficient and nondegenerate Hopf theorem
% DEPENDS ON: 00_canonical_model.tex; 01_notation_and_assumptions.tex; shared_macros.tex; P11, P12. Optional dependencies are stated below.
% ASSUMPTIONS: P11--P12 hypotheses; smooth DM field near the equilibrium; explicitly reconciled Hermitian pairing, adjoint vector and Taylor factorial conventions.
% TARGET STATEMENTS: Registry C077, C078, C079, C080, C081, C082, C083, C148. Bilinear and trilinear derivative tensors | Fourier selection rules and resolvent domains | First Lyapunov coefficient and rigorous sign or nonvanishing | Assembled nondegenerate Hopf theorem and branch stability conditions
% REQUIRED CHECKS: Audit the blueprint adjoint eigenvalue sign against <p,v>=p^*v; tensor factors; eigenvector scaling; modular aliasing of 2m; inverse/resolvent existence; software sign conventions; transverse stability.
% POSSIBLE FAILURE MODES: Copying an inconsistent adjoint convention; l1 value changes with untracked normalization; a sign claim based solely on floating point; calling a branch stable from l1 alone.
% LITERATURE DEPENDENCIES: Later: primary normal-form reference and independent software documentation for the exact l1 convention.
% NUMERICAL CROSS-CHECKS: N01: two independent l1 implementations plus derivative-tensor checks; N02: periodic-branch direction and Floquet multipliers.
% TODO: B07: resolve adjoint/Taylor/l1 conventions; B06: choose a rigorous sign-certification strategy. Later derive each statement independently, write the proof, verify it, and commission a separate audit. No proof or computation is supplied at initialization.

\section{P13: First Lyapunov coefficient and nondegenerate Hopf theorem}
\label{sec:P13}

\subsection{TODO: Bilinear and trilinear derivative tensors}
\label{slot:P13-1}

\subsection{TODO: Fourier selection rules and resolvent domains}
\label{slot:P13-2}

\subsection{TODO: First Lyapunov coefficient and rigorous sign or nonvanishing}
\label{slot:P13-3}

\subsection{TODO: Assembled nondegenerate Hopf theorem and branch stability conditions}
\label{slot:P13-4}
